%% =====================================================================
%% PART 1 WRITING SKELETON — empty sections with instructions
%% Paste over the Part 1 block in your Overleaf main.tex.
%% Write your own text where each "WRITE:" comment sits, then DELETE the
%% comments before you submit. Nothing here is content — only prompts.
%%
%% Total: 750–1000 words excluding references. No code or pseudocode.
%% 5–10 additional peer-reviewed in-text citations beyond the paper itself.
%% Paper: \citep{trippas2025adapting}  <- use your own bib key
%% =====================================================================

\section{Part 1: Executive Summary}
\label{sec:part1}

%% ---------------------------------------------------------------------
\subsection{Purpose and Audience}
%% TARGET: ~120 words. 1–2 citations.
%% Answer, in this order:
%%   1. What is the paper, who wrote it, and what kind of piece is it
%%      (position chapter, not an experimental study)?
%%   2. What problem or gap prompted it? (from notes heading 1)
%%   3. What does it set out to do? (notes heading 2)
%%   4. Who is it written for, and what is your evidence for that?
%%      Look at where it was published, the terminology, and what it asks
%%      the reader to do.
%% CHECK: a reader who has never seen the paper should now know why it exists.

% WRITE: purpose and audience here.

%% ---------------------------------------------------------------------
\subsection{Main Argument and Contributions}
%% TARGET: ~220 words. 2–3 citations.
%% Answer:
%%   1. What is the paper's central claim, in one sentence of your own?
%%   2. What does it actually deliver — a framework, a synthesis, a research
%%      agenda, a set of definitions? Name them (notes heading 3).
%%   3. Which 2–3 contributions matter most for a data science practitioner,
%%      and why those?
%%   4. One sentence positioning it against other work: does it extend,
%%      challenge or synthesise what others argue? This is where an external
%%      citation earns its place.
%% AVOID: walking through the chapter section by section. Summarise the
%%        argument, don't narrate the table of contents.

% WRITE: main argument and contributions here.

%% ---------------------------------------------------------------------
\subsection{Case Study}
%% TARGET: ~250 words. 1–2 citations (you need a real source for the case).
%% CONSTRAINT: must NOT be one of the paper's own scenarios (its Section 5:
%%   work, knowledge-base access via conversational agents, learning and
%%   teaching, research, personal information management).
%% Answer:
%%   1. What is the case? Who are the users, what data is involved, and what
%%      decision or task does the system support? Be concrete.
%%   2. Why did you choose it — which of the paper's ideas does it test?
%%   3. One sentence saying explicitly how it differs from the paper's own
%%      scenarios. (Shows the marker you respected the constraint.)
%%   4. What happens when the paper's framing is applied to it?
%% CHECK: a marker should be able to picture the system and its users.

% WRITE: case study here.

%% ---------------------------------------------------------------------
\subsection{Challenges}
%% TARGET: ~180 words. 1–2 citations.
%% Pick the 3–4 challenges from your notes (heading 5) that your case really
%% shows. For EACH one:
%%   - name the challenge in the paper's terms,
%%   - show how it appears concretely in your case,
%%   - say what it costs (a wrong decision, a harmed group, wasted spend).
%% AVOID: listing every challenge the paper mentions. Depth beats coverage.

% WRITE: challenges here.

%% ---------------------------------------------------------------------
\subsection{Opportunities}
%% TARGET: ~150 words. 1–2 citations.
%% For 2–3 opportunities from your notes (heading 6):
%%   - what becomes possible in your case that wasn't before,
%%   - what would have to be true for that benefit to be real
%%     (data, evaluation, governance, human oversight),
%%   - who benefits, and who might not.
%% Balance matters: an opportunity with a condition attached reads as
%% analysis; an unconditional one reads as marketing.

% WRITE: opportunities here.

%% ---------------------------------------------------------------------
\subsection{Conclusion}
%% TARGET: ~80 words. No new citations.
%% Answer:
%%   1. What is the paper's value for data science practice, in your judgement?
%%   2. What does your case add to, or qualify in, its argument?
%% AVOID: repeating the summary. One evaluative point is enough.

% WRITE: conclusion here.

%% =====================================================================
%% BEFORE YOU MOVE ON — checklist
%%  [ ] 750–1000 words excluding references (paste into a word counter)
%%  [ ] 5–10 peer-reviewed citations beyond the paper, all resolving in the PDF
%%      (Task 1 feedback: journal/conference sources, not just websites)
%%  [ ] Case study is not one of the paper's own scenarios, and you said so
%%  [ ] Motivation, aims and contributions all explicitly covered
%%  [ ] No code or pseudocode
%%  [ ] Every "WRITE:" and instruction comment deleted
%% =====================================================================
